\section{Normalization and the comparison with the raw calculus}
\label{sec:fullnf-metatheory}

This section supplies the metatheory of the computation built into the
new syntax.  Write $\mathcal N$ for the raw native calculus with its two
computation axioms, and $\mathcal F$ for the six-rule normalized calculus.
All statements about normalization are external statements about finite
syntax; they are not new axioms of the numerical theory.

\subsection{Strong normalization and unique normal forms}

\begin{definition}[Contextual computation]
On raw terms, close the two reductions
\begin{equation}
\begin{aligned}
 I(0,a;u.h)&\longrightarrow a,\\
 I(Sn,a;u.h)&\longrightarrow h[I(n,a;u.h)/u]
\end{aligned}
\label{eq:fullnf-reductions}
\end{equation}
under every term context, including iterator-update contexts.  Substitution
avoids capture.  A term is \emph{strongly normalizing} if it admits no
infinite reduction sequence.  For such a term $t$, let $\nu(t)$ be its
maximum reduction length, and let $|t|$ be its syntax-tree size.
\end{definition}

The maximum exists: a finite term has only finitely many one-step reducts,
and a finitely branching well-founded reduction tree has finite height.

\begin{lemma}[Iterator closure]
\label{lem:fullnf-closure}
Suppose $n$ and $a$ are strongly normalizing, and $h[v/u]$ is strongly
normalizing for every strongly normalizing term $v$.  Then
$I(n,a;u.h)$ is strongly normalizing.
\end{lemma}
\begin{proof}
The body $h$ itself is strongly normalizing, by taking $v=u$.
Use lexicographic induction on
\[
  \bigl(\nu(n),\ |n|,\ \nu(a)+\nu(h)\bigr).
\]
It suffices to show that every one-step reduct of the iterator is strongly
normalizing.  A reduction in $n$ decreases the first coordinate.  A reduction
in $a$ or $h$ decreases the last coordinate.  In the latter case, if
$h\longrightarrow h'$, then $h'[v/u]$ is a reduct of $h[v/u]$ for every
$v$, so the required universal hypothesis persists.  Here we use the
immediate fact that contextual computation is stable under capture-free
substitution.

The zero root reduct is $a$.  At a successor root, write $n=Sp$.
Then $\nu(p)=\nu(n)$ and $|p|<|n|$.  The induction hypothesis therefore
makes $I(p,a;u.h)$ strongly normalizing.  The hypothesis on $h$ makes
$h[I(p,a;u.h)/u]$ strongly normalizing, as required.
\end{proof}

\begin{theorem}[Strong normalization]
\label{thm:fullnf-sn}
Substitution of strongly normalizing terms for the free variables of any
finite raw term produces a strongly normalizing term.  In particular,
every finite raw term is strongly normalizing.
\end{theorem}
\begin{proof}
Induct on the original term.  The variable, zero, and successor cases are
immediate.  At an iterator, the induction hypotheses for its counter and
initializer give strongly normalizing instances.  The body induction
hypothesis applies to \emph{every} additional assignment $u:=v$ with $v$
strongly normalizing.  It supplies exactly the universal premise of
Lemma~\ref{lem:fullnf-closure}.  Freshen the bound variable before making
these assignments.  Finally substitute each free variable with itself.
\end{proof}

\begin{lemma}[Confluence and canonical syntax]
\label{lem:fullnf-confluence}
Every raw term $t$ has a unique normal form $\operatorname{nf}(t)$, up to
renaming of bound variables, or literally when de Bruijn indices are used.
Its iterator subterms have normal children and counters headed by neither
$0$ nor $S$.
\end{lemma}
\begin{proof}
The reduction is locally confluent.  Disjoint redexes commute.  The two
root patterns in~\eqref{eq:fullnf-reductions} are disjoint.  For a redex
nested inside a root contraction, contract all its residual copies after
the root step.  The zero rule may erase the update.  The successor rule
may copy the update and recursive iterator into several occurrences of
the bound variable, but there are only finitely many copies; substitution
stability gives a common reduct.  These exhaust the overlaps.

Strong normalization and local confluence give confluence by induction on
reduction height.  A term is normal precisely when its children are normal
and neither root iterator rule applies, giving the stated grammar.
\end{proof}

\subsection{Hereditary substitution and the shared walk}

For a substitution $\sigma$, write
\[
 t\langle\!\langle\sigma\rangle\!\rangle
       :=\operatorname{nf}(t[\sigma]).
\]
Formula normalization preserves the constructors $=$, $Q$, and $D$, and
normalizes all their numerical terms, including terms under binders.

\begin{lemma}[Substitution coherence]
\label{lem:fullnf-substitution}
For every raw term $t$ and substitution $\sigma$,
\begin{equation}
 \operatorname{nf}(t[\sigma])
 =\operatorname{nf}\bigl(\operatorname{nf}(t)
                     [x\mapsto\operatorname{nf}(\sigma x)]\bigr).
\label{eq:fullnf-coherence}
\end{equation}
Consequently normalized substitution has the identity, composition,
weakening, and binder-lifting laws.  The same assertions hold for formulas.
\end{lemma}
\begin{proof}
The reduction $t\longrightarrow^*\operatorname{nf}(t)$ survives raw
substitution.  Reducing each of the finitely many substituted copies of
$\sigma x$ to its normal form gives conversion to the right-hand input
in~\eqref{eq:fullnf-coherence}.  Uniqueness of normal forms proves the
equation.  Apply it twice and use associativity of raw substitution to get
\[
 (t\langle\!\langle\sigma\rangle\!\rangle)
       \langle\!\langle\tau\rangle\!\rangle
 =t\langle\!\langle
          x\mapsto(\sigma x)\langle\!\langle\tau\rangle\!\rangle
       \rangle\!\rangle.
\]
Renamings and lifted substitutions are special cases.  Formula constructors
have no computation rules, so structural recursion gives the formula claims.
\end{proof}

The executable walk first processes the children of a node, then contracts
an exposed iterator redex.  At a successor it passes the raw recursive
iterator into hereditary substitution; when the bound variable actually
occurs, the substitution's shifting walk normalizes that replacement.
If the update ignores its variable, the recursive iterator is erased
without evaluation.  Thus every computational step contracts a genuine
redex in the materialized raw substituted term.  Between contractions the
algorithm traverses finite subtrees and renames indices.  This is ordinary
normalization with raw substitution fused into traversal: strong
normalization proves termination, and the normal, reachable result is
$\operatorname{nf}(t)$ by uniqueness.  The same argument identifies the
fused substitution with $\operatorname{nf}(t[\sigma])$.

\subsection{Equality inside an update binder}

There is a small but essential proof obligation when comparing the new
syntax with $\mathcal N$.  Capture-free one-term transport alone does not
allow a formerly free variable in its equality premise to become bound
inside an iterator.  The missing congruence is derived using equation
induction, which $\mathcal N$ already proves.

\begin{lemma}[Iterator-body extensionality]
\label{lem:fullnf-body}
If $u$ is fresh for $\Gamma$ and
$\Gamma\vdash_{\mathcal N}h(u)=k(u)$, then, for all terms $n,a$,
\[
 \Gamma\vdash_{\mathcal N}I(n,a;u.h)=I(n,a;u.k).
\]
\end{lemma}
\begin{proof}
Choose $z$ fresh for all passive parameters and induct on the literal equation
\[
 E(z):=I(z,a;u.h)=I(z,a;u.k).
\]
Both zero computations give the base.  Assuming $E(z)$, successor
computation, numerical substitution in the pointwise equality, and ordinary
equality transport give
\[
\begin{aligned}
 I(Sz,a;h)&=h(I(z,a;h))\\
          &=k(I(z,a;h))\\
          &=k(I(z,a;k))\\
          &=I(Sz,a;k).
\end{aligned}
\]
Equation induction gives $E(n)$.  No induction on a quantified formula
or on an implication is used.
\end{proof}

\begin{lemma}[Normalization is provable in the raw calculus]
\label{lem:fullnf-provable}
For each fixed finite raw term $t$,
\[
 \vdash_{\mathcal N}t=\operatorname{nf}(t).
\]
For each raw formula $A$, the formulas $A$ and $\operatorname{nf}(A)$ are
interderivable in $\mathcal N$, uniformly under additional retained contexts.
\end{lemma}
\begin{proof}
Each root reduction is an old computation axiom.  Transport handles
replacement in a successor, counter, or initializer; repeated use of
Lemma~\ref{lem:fullnf-body} handles replacement under update binders.
Thus each step of a finite normalization derivation is a provable equality.
Transitivity gives the term assertion.

Induct externally on formula syntax.  The equation case uses the two term
equalities and ordinary equality reasoning.  At a $Q$ node, convert its
new zero instance to the old zero instance, specialize the old closure at
$Sz$, convert the result, and apply the derived case introduction for $Q$;
reverse the conversions for the converse.  At a $D$ node, eliminate the
source at a fresh witness, perform the preceding closure conversion, and
introduce the target at that same witness.  These constructions are uniform
under passive contexts and avoid capture.  They are binder congruence,
not a claim that an open numerical variable may be captured by transport.
\end{proof}

\subsection{Both consequence translations and both composites}

Let $T$ normalize every numerical term of a raw formula, and let $U$ regard
a canonical formula as raw syntax.  Let $\mathcal F^+$ be the auxiliary
normalized calculus with separate general case introduction and direct
equation-progress introduction.  The compound-envelope argument gives
the same derivable sequents in $\mathcal F^+$ and $\mathcal F$.

\begin{theorem}[Exact comparison]
\label{thm:fullnf-comparison}
For every raw sequent and every canonical sequent,
\begin{align*}
 \Gamma\vdash_{\mathcal N}A
   &\quad\Longleftrightarrow\quad
       T\Gamma\vdash_{\mathcal F}TA,\\
 \Delta\vdash_{\mathcal F}B
   &\quad\Longleftrightarrow\quad
       U\Delta\vdash_{\mathcal N}UB.
\end{align*}
Moreover $TU(B)=B$ literally, and $A$ and $UT(A)$ are uniformly
interderivable in $\mathcal N$.
\end{theorem}
\begin{proof}
First translate a raw derivation into $\mathcal F^+$ by induction on the
derivation.  Both old computation equations normalize to reflexivity.
The remaining rule instances normalize to their corresponding instances
by Lemma~\ref{lem:fullnf-substitution}; literal-equation status and
eigenvariable conditions are preserved.  Replace auxiliary introductions
by their six-rule derivations.  This avoids the invalid direct argument
that a nonzero raw witness must remain nonzero after normalization.

Conversely induct on an $\mathcal F$ derivation.  Reflexivity is already
derived in $\mathcal N$.  Before applying an old rule, use
Lemma~\ref{lem:fullnf-provable} to convert its normalized instance premises
to raw instances; convert its conclusion back afterwards.  A case-mode
witness introduction uses old case introduction.  A progress-mode
introduction first uses old equation progress on the instantiated literal
equation, then obtains the witnessed $D$ formula by old case introduction
and closure specialization.  Hypothesis, transport, $Q$ elimination, and
$D$ elimination translate in the same way.  Context conversions use the
admissible assumption cut and weakening of the old calculus.

The composite $TU$ is the identity because canonical terms are normal.
The other composite is formula normalization, handled by
Lemma~\ref{lem:fullnf-provable}.  Combining these facts with the two forward
translations proves both displayed equivalences, including the directions
whose context was normalized.
\end{proof}

These maps interpret the term presentations modulo provable equality;
normalization is not an additional numerical function symbol.  Composing
them with the inherited native-to-iPRA bi-interpretation gives the required
formula round trips and equiconsistency.  The numerical domains remain
whole, equality remains actual equality, and both domain isomorphisms have
graph $x=y$.  Only finite \emph{term syntax} is identified modulo computation;
there is no quotient or restriction of the numerical domain.

\subsection{The uniform cost is external}

Every fixed scalar term denotes a primitive-recursive function, but the
normalizer on \emph{all} finite term codes is not primitive recursive.
Otherwise normalizing one-variable scalar terms on numeral inputs would
provide a primitive-recursive evaluator $\operatorname{Eval}(e,n)$ for
all primitive-recursive functions.  Primitive-recursive syntactic coding
and extraction of a normal numeral's value would then make
$g(e)=\operatorname{Eval}(e,e)+1$ primitive recursive, contradicting its
own scalar-term index.  Invalid codes can be assigned the default value
zero, so they do not affect the diagonal argument.

The size comparison therefore counts a total recursive normalizing checker,
with potentially enormous running time and output.  It does not assert
that iPRA proves a uniform normalization theorem.  If one separately
requires a primitive-recursive certificate relation, finite reduction
traces may be supplied and checked locally; that is a different checker
whose added code must also be counted.  The new normalization and comparison
results above are written mathematical proofs, not claims of a new
proof-assistant verification.
